\section{Spray Physics \& Atomization}

\subsection{Introduction}
This chapter describes the physics of propellant atomization, spray formation, and evaporation. The models track the transformation of liquid propellant streams into droplets and their subsequent gasification.

\subsection{Nomenclature}
\begin{description}
    \item[\(J\)] Momentum flux ratio [dimensionless]
    \item[\(TMR\)] Thrust momentum ratio [dimensionless]
    \item[\(\theta\)] Spray angle [\si{rad}]
    \item[\(We\)] Weber number [dimensionless]
    \item[\(Oh\)] Ohnesorge number [dimensionless]
    \item[\(D_{32}\)] Sauter Mean Diameter (SMD) [\si{m}]
    \item[\(\tau_{evap}\)] Evaporation time [\si{s}]
    \item[\(x^*\)] Evaporation length [\si{m}]
    \item[\(\sigma\)] Surface tension [\si{N/m}]
\end{description}

\subsection{Momentum and Spray Angle}

\subsubsection{Momentum Flux Ratio}
The momentum flux ratio \(J\) characterizes the relative strength of the oxidizer and fuel streams:
\begin{equation}
J = \frac{\rho_O u_O^2}{\rho_F u_F^2}
\end{equation}
\coderef{engine/core/spray.py}{31}

\subsubsection{Thrust Momentum Ratio}
The Thrust Momentum Ratio (TMR) is calculated as:
\begin{equation}
TMR = \frac{J}{1 + MR} \quad (\text{if } MR > 0, \text{ else } J)
\end{equation}
\coderef{engine/core/spray.py}{55}

\subsubsection{Spray Angle}
The spray angle is determined using the half-angle tangent form:
\begin{equation}
\tan(\theta/2) = k J^n
\end{equation}
\begin{equation}
\theta = 2 \arctan(k J^n)
\end{equation}
\coderef{engine/core/spray.py}{120-121}

Alternatively, using TMR:
\begin{equation}
\theta = \arccos\left(\frac{1}{1 + TMR^{0.75}}\right)
\end{equation}
\coderef{engine/core/spray.py}{146-148}

\subsection{Atomization Quality}

\subsubsection{Weber Number}
The Weber number represents the ratio of inertial forces to surface tension:
\begin{equation}
We = \frac{\rho u^2 d_{char}}{\sigma}
\end{equation}
\coderef{engine/core/spray.py}{177}

\subsubsection{Ohnesorge Number}
The Ohnesorge number relates viscous forces to surface tension and inertial forces:
\begin{equation}
Oh = \frac{\mu}{\sqrt{\rho \sigma d_{or}}}
\end{equation}
\coderef{engine/core/spray.py}{207}

\subsubsection{Sauter Mean Diameter (SMD)}
For general injectors, the Lefebvre correlation is used:
\begin{equation}
D_{32} = C d_{or} We^{-m} Oh^p
\end{equation}
\coderef{engine/core/spray.py}{246}

For pintle injectors, a physics-based correlation using relative velocity is applied:
\begin{equation}
D_{32} = C L_{open} We_{rel}^{-n} (1 + B Oh_f)^p
\end{equation}
where \(We_{rel}\) is based on \(V_{rel} = \sqrt{u_O^2 + u_F^2}\).
\coderef{engine/core/spray.py}{405-417}

\subsection{Evaporation and Penetration}

\subsubsection{Evaporation Time}
The evaporation time follows the \(D^2\)-law:
\begin{equation}
\tau_{evap} = K D_{32}^2
\end{equation}
\coderef{engine/core/spray.py}{267}

\subsubsection{Evaporation Length}
The characteristic evaporation length \(x^*\) is:
\begin{equation}
x^* = U_{rel} \tau_{evap}
\end{equation}
\coderef{engine/core/spray.py}{288}

If turbulence corrections are enabled:
\begin{equation}
x^* = x^* \times \text{clamp}\left(\frac{1}{1 + G_{turb} I_{mix}}, 0.3, 1.0\right)
\end{equation}
\coderef{engine/core/injectors/pintle.py}{267-269}

